% Currículum actualizado — Juan Esteban Cardona Blandón
\documentclass[11pt,a4paper,sans]{moderncv}

% Estilo y color
\moderncvstyle{banking}
\moderncvcolor{blue}

% Márgenes
\usepackage[scale=0.8]{geometry}
\usepackage{graphicx}
\usepackage{import}
\usepackage{booktabs}
\usepackage{xcolor}
\usepackage{tabularx}
\usepackage[export]{adjustbox}
\usepackage[hidelinks]{hyperref}
\usepackage[spanish]{babel}

\hypersetup{
    colorlinks=true,
    linkcolor=blue,
    urlcolor=blue,
    citecolor=blue
}
\setlength{\footskip}{136.00005pt}

% Datos personales
\name{Juan Esteban}{Cardona Blandón}
\title{Hoja de Vida}
\photo[64pt][0.6pt]{photo}
\address{Pereira, Colombia}{}{}
\email{juan.cardona13@utp.edu.co}
\extrainfo{+57 3012213085 \textbullet{} C.C. 1113858851 \textbullet{} Fecha de nacimiento: 3 de septiembre de 2004}

\quote{Estudiante de Ingeniería de Sistemas en la Universidad Tecnológica de Pereira, con experiencia práctica en desarrollo de software (Python, JavaScript/TypeScript, React, Tailwind), despliegue y orquestación de contenedores (Docker, Kubernetes) y diseño digital con lenguajes HDL para FPGA. Interesado en prácticas profesionales donde pueda aportar en proyectos reales, aprender metodologías ágiles y continuar fortaleciendo habilidades técnicas y de trabajo en equipo.}

\makeatletter
\@ifpackageloaded{moderncvstylebanking}{%
\let\oldmakecvtitle\makecvtitle
\renewcommand*{\makecvtitle}{%
  {\centering\framebox{\includegraphics[width=\@photowidth]{\@photo}}\par\vspace{12pt}}%
  \oldmakecvtitle%
}%
}{%
}
\makeatother

\microtypesetup{expansion=false,protrusion=false}

\begin{document}

\makecvtitle


\section{Educación}
\cventry{2022--2025}{Ingeniería de Sistemas}{Universidad Tecnológica de Pereira}{}{}{8º semestre en curso. Formación en algoritmos, arquitectura de computadoras, bases de datos, redes y sistemas embebidos.}


\section{Experiencia Académica}
\cventry{2025}{Monitor Académico}{Universidad Tecnológica de Pereira}{}{}{Apoyo en asignaturas de programación y arquitectura de computadores: resolución de dudas, desarrollo de material auxiliar y acompañamiento a estudiantes.}


\section{Habilidades Técnicas}
\cvitem{Lenguajes}{Python, JavaScript, TypeScript, C, C++}
\cvitem{Frameworks / Frontend}{Django, React, Tailwind CSS}
\cvitem{DevOps / Contenedores}{Docker, Docker Compose, Kubernetes}
\cvitem{Bases de datos}{SQLite, MySQL, PostgreSQL (básico)}
\cvitem{Hardware / HDL}{Verilog, VHDL, FPGA (DE1-SoC), diseño digital}
\cvitem{Herramientas}{Git, GitHub, Linux}


\section{Proyectos Relevantes}
\cvitem{GitHub}{\url{https://github.com/Juanes7222} — Repositorio con proyectos personales y académicos.}

\cvitem{MapReduce}{Implementación y ejercicios sobre el modelo MapReduce. \\ \url{https://github.com/Juanes7222/MapReduce}}

\cvitem{Kubernetes-SD}{Prácticas y experimentos con Kubernetes; incluye manifiestos y ejemplos de despliegue. \\ \url{https://github.com/Juanes7222/Kubernetes-SD}}

\cvitem{Otros proyectos}{Procesamiento de imágenes en Python (FFT), aplicaciones web con Django, diseño digital en FPGA con Verilog/VHDL.}

\cvitem{Contribuciones}{RiscVSiriusStudio — Contribuciones al proyecto de arquitectura y herramientas para RISC-V. \\ \url{https://github.com/LabSirius/RiscVSiriusStudio}}


\section{Formación Complementaria}
\cventry{2025}{Diplomado en Responsabilidad Social}{Universidad Tecnológica de Pereira}{}{}{Curso y actividades relacionadas con responsabilidad social universitaria.}

\cventry{2024--2025}{Mentorías de habilidades blandas}{Banco de Bogotá}{Programa de beca}{}{Participación en talleres de liderazgo, comunicación asertiva y trabajo en equipo.}


\section{Habilidades Blandas}
\cvitem{Pensamiento analítico}{Alta capacidad para verificar, revisar y asegurar la calidad del trabajo con enfoque riguroso.}
\cvitem{Autocontrol}{Mantengo estabilidad emocional y orden incluso bajo presión.}
\cvitem{Comunicación}{Explico ideas técnicas de forma clara, respetuosa y estructurada.}
\cvitem{Trabajo en equipo}{Colaboración armónica, orientada a soluciones prácticas y eficientes.}
\cvitem{Responsabilidad}{Cumplimiento de compromisos, confiabilidad y atención al detalle.}
\cvitem{Organización}{Planificación, priorización y ejecución disciplinada de tareas.}
\cvitem{Prudencia}{Toma de decisiones consciente, evaluando riesgos y escenarios posibles.}
\cvitem{Aprendizaje}{Adaptabilidad a nuevas tecnologías y mejora continua.}
\cvitem{Pruebas DISC}{\href{https://drive.google.com/file/d/1jkTzgjt8AS6A15DcztzRoxmq90OAagMo/view?usp=sharing}{Ver resultados completos}}

\section{Idiomas}
\cvitem{Español}{Nativo}
\cvitem{Inglés}{Intermedio (lectura de documentación técnica y comunicación básica)}

\end{document}